\section{A shared CRT compiler with binomial bounds}
\label{sec:shared-crt}

All variables in this section range over $\mathbb N=\{0,1,2,\ldots\}$.
Let $E(k,z,w)$ and $F(k,z,w)$ be fixed polynomials with nonnegative
integer coefficients, where $k$ is one index, $z$ is a finite tuple of
external parameters, and $w=(w_1,\ldots,w_d)$ is a finite witness tuple.
The nonnegative coefficient hypothesis is used only for monotonicity.
For a proposed vector of witness bounds $u=(u_1,\ldots,u_d)$, choose a
natural $V$ with $u_j\le V$ for every $j$, and put
\[
 H=E(N,z,u)+F(N,z,u)+N+V+1,\qquad T=H!,\qquad
 B=\prod_{k=0}^{N-1}(1+(k+1)T).
\]
The empty product is $1$.  In particular $T\ge1$, $B\ge1$, and
$H\ge u_j+1$ for every coordinate.

\begin{theorem}[Shared binomial certificate]
\label{thm:shared-binomial}
For fixed $N,z,u,V$ as above, the following assertions are equivalent:
\begin{enumerate}
\item For every $k<N$ there is $w\in\mathbb N^d$ with
      $w_j\le u_j$ for all $j$ and $E(k,z,w)=F(k,z,w)$.
\item There are $C,A_1,\ldots,A_d\in\mathbb N$ such that
\begin{align}
 B&\mid\binom{A_j}{u_j+1} &&(1\le j\le d),\label{eq:choose-bounds}\\
 B&\mid1+(C+1)T,\label{eq:index-divisibility}\\
 E(C,z,A)&\equiv F(C,z,A)\pmod B.\label{eq:evaluation-congruence}
\end{align}
\end{enumerate}
\end{theorem}

\begin{proof}
Write $D_k=1+(k+1)T$.  We have $\gcd(D_k,T)=1$.
The numbers $D_0,\ldots,D_{N-1}$ are pairwise coprime: if $i<j<N$
and $g=\gcd(D_i,D_j)$, then $g\mid(j-i)T$ and $\gcd(g,T)=1$,
so $g\mid j-i$.  But $j-i\le N\le H$, hence $j-i\mid H!=T$,
and $g=1$.  Consequently $\gcd(B,T)=1$ as well.

Suppose first that the arithmetic certificate exists.  Fix $k<N$.
Choose a prime $p\mid D_k$; this is possible since $D_k>1$.
Then $p\mid B$ and $p\nmid T$.  A prime at most $H$ divides $H!$,
so $p>H$.  For every $j$, the integer identity
\[
 (u_j+1)!\binom{A_j}{u_j+1}
       =\prod_{r=0}^{u_j}(A_j-r)
\]
and \eqref{eq:choose-bounds} show that $p$ divides the product on the
right.  This identity also holds when $A_j\le u_j$, because both sides
are zero.  Primality therefore supplies $w_j\in\{0,\ldots,u_j\}$
with $A_j\equiv w_j\pmod p$.  Comparing
\eqref{eq:index-divisibility} with $D_k\equiv0\pmod p$, and cancelling
$T$ modulo $p$, gives $C\equiv k\pmod p$.  Polynomial evaluation and
\eqref{eq:evaluation-congruence} imply
\[
 E(k,z,w)\equiv F(k,z,w)\pmod p.
\]
Monotonicity gives
\[
 0\le E(k,z,w)\le E(N,z,u)<p,
 \qquad 0\le F(k,z,w)\le F(N,z,u)<p.
\]
These congruent natural numbers are equal, proving the bounded row.

Conversely, choose one bounded witness $w^{(k)}$ for each $k<N$.
By the Chinese remainder theorem and the pairwise coprimality just proved,
choose natural $C,A_j$ satisfying
\[
 C\equiv k\pmod{D_k},\qquad
 A_j\equiv w_j^{(k)}\pmod{D_k}
 \quad(k<N,\ 1\le j\le d).
\]
The latter congruence makes the falling-factorial product vanish modulo
$D_k$, since its factor indexed by $w_j^{(k)}$ vanishes.  Thus
\[
 B\mid (u_j+1)!\binom{A_j}{u_j+1}.
\]
Now $(u_j+1)!$ divides $H!=T$ and $\gcd(B,T)=1$; we may cancel this
factor to obtain \eqref{eq:choose-bounds}.  The index congruences give
\eqref{eq:index-divisibility}; the valid rows and polynomial congruence
give \eqref{eq:evaluation-congruence}.  Pairwise coprimality combines
the congruences modulo the $D_k$ into congruences modulo $B$.
When $N=0$, there are no row obligations and $B=1$, so both directions
hold without any separate side condition.
\end{proof}

\paragraph{Why the factorial cancellation matters.}
The repository's supplied CRT certificate uses
$B\mid(u_j+1)!\binom{A_j}{u_j+1}$ separately for every coordinate.
Its exact source does not require $H\ge u_j+1$, so simply dropping these
factorials from that source would be unjustified.  The explicit added bound
$V+1$ repairs this issue.  It also permits heterogeneous bounds at no
additional factorial cost.  A representation of each binomial coefficient
is still necessary; the theorem does not eliminate that obligation.

\subsection{A finite arithmetic expression for the modulus product}

The product defining $B$ is explanatory notation and is eliminated before
the final polynomial is issued.  Set
\[
 Q=(1+(N+1)T)^N+2,\qquad L=TQ-1.
\]
Then the exact natural-number identity is
\begin{equation}
 \label{eq:product-expression}
 B=\left(T^N N!\binom{Q+N}{N}\right)\bmod L.
\end{equation}
Indeed, the numerator equals
$\prod_{i=1}^N(TQ+iT)$ and reduces modulo $L$ to
$\prod_{i=1}^N(1+iT)=B$.  Further,
\[
 B\le(1+NT)^N<(1+(N+1)T)^N+1\le TQ-1=L,
\]
with the strict inequality also valid for $N=0$, when the left side is
$1$ and the middle upper bound is $2$.  Thus $B$ is the remainder.
The denominator is positive because $T\ge1$ and $Q\ge3$.
This is the same explicit product mechanism used by the repository's
\nolinkurl{BoundedForallRepresentation.productGraph}; the elementary
calculation above records the needed identity independently.

Powers, binomial coefficients, factorials, and the final remainder are each
replaced by their fixed finite polynomial graph.  In particular, no
variable-length product, variable exponent, or bounded universal quantifier
remains in the emitted ordinary polynomial.

\subsection{Keeping the majorants as arithmetic circuits}

A signed arithmetic circuit $P$ admits natural polynomial circuits $E,F$
such that $P=E-F$.  The recursive pairs are $(X,0)$ for a variable,
$(c,0)$ or $(0,-c)$ for an integer constant according to its sign,
\[
 (E_1,F_1)+(E_2,F_2)=(E_1+E_2,F_1+F_2),
\]
\[
 (E_1,F_1)-(E_2,F_2)=(E_1+F_2,F_1+E_2),
\]
and
\[
 (E_1,F_1)(E_2,F_2)=
 (E_1E_2+F_1F_2,\ E_1F_2+F_1E_2).
\]
Every operation preserves nonnegative coefficients.  Repeated subexpressions
are shared as nodes in a finite directed acyclic graph.  Expanding this DAG
into a tree is mathematically harmless but can make the representation far
larger.  The compiler uses node values as natural witnesses and emits their
quadratic defining equations only once per use-site evaluation.

\section{Eliminating local Horn wires before bounded compilation}

The pinned \nolinkurl{affine_shared_step.json} source contains twelve Horn
branches.  Its existential coordinates include thirteen shared names
$q_0,\ldots,q_{12}$.  In each branch the actual wire definitions are
triangular polynomials with nonnegative coefficients, and the first rows of
that branch assert exactly $q_i-f_i=0$.  Consequently, on natural inputs,
the definitions determine a unique natural value for every wire used by that
branch.  Substitution removes all thirteen existential wire coordinates.

There are fifty-six branch-specific definition occurrences and thirty-four
remaining branch residuals.  These are verified from the pinned JSON by exact
sparse-polynomial comparison.  Branch definitions use separate scopes.  This
does not assume one set of wire values makes all branches vanish: existential
quantification distributes over the finite disjunction, and precisely one
vanishing branch is sufficient.  The product of the twelve branch sums of
squares implements that disjunction.

The affine step originally has twenty-three natural witness coordinates.
Wire elimination leaves ten:
\[
 p,q,v,a,b,c,d,\delta_p,\delta_q,\delta_v.
\]
Here $(p,q)\to v$ is the rule, $a,b,c,d$ are its four schema witnesses,
and the three deltas implement $p,q,v\le m$.
Using Theorem~\ref{thm:shared-binomial} with heterogeneous bounds removes
the latter three witnesses too: bound $p,q,v$ directly by $m$ in the CRT
certificate.  Bound all remaining per-row witnesses by $V=m+U$, where
$U$ is one additional outer natural variable.  Such $U$ always exists for
any finite valid certificate.

\subsection{Thirteen per-row coordinates}

Retain the five outer parameters $t,x,y,m,M_0$ and put $M=M_0+1$.
The carrier condition is $M=(m+1)^t$.  At row $i<t$, set
$D_i=1+(i+1)y$ as an arithmetic expression, without introducing a
separately bounded row coordinate.  Besides the seven Horn coordinates
$p,q,v,a,b,c,d$, use six row coordinates $S,S',u,r,u',r'$.
The four beta equations are
\begin{align*}
 x-S-D_i u&=0,& S+r+1-D_i&=0,\\
 x-S'-(D_i+y)u'&=0,& S'+r'+1-(D_i+y)&=0.
\end{align*}
Since $D_i,D_i+y\ge1$, these say exactly that $S$ and $S'$ are
the consecutive beta digits.  The remaining accumulator equation is
\[
 S'(1+vM)=S(1+pM)(1+qM).
\]
Let $P$ be the product of the twelve substituted Horn branch sums of squares,
plus the square of this accumulator residual and the squares of the four
beta residuals.  It is a fixed nonnegative polynomial in the index,
the external parameters and these thirteen row coordinates.  Split its
arithmetic circuit as $P=E-F$ as above.

For the first three encoded coordinates use the bound $m$, and for the
other ten use $V=m+U$.  Form
\[
 H=E(t,z,u)+F(t,z,u)+t+V+1,
 \qquad u=(m,m,m,V,\ldots,V),
\]
where $z=(t,x,y,m,M_0)$ denotes the external parameters (the polynomial
need not actually use $t$ except through the index bound).
The shared certificate then expresses all rows with thirteen binomial
graphs and no per-coordinate factorial graphs.

The endpoints need four natural witnesses $u_0,r_0,u_t,r_t$ and the rows
\begin{align*}
 x-(M_0+2)-(y+1)u_0&=0,\\
 M_0+r_0+2-y&=0,\\
 x-1-(1+(t+1)y)u_t&=0,\\
 r_t+2-(1+(t+1)y)&=0.
\end{align*}
They assert $\beta(0)=M+1$ and $\beta(t)=1$ exactly.

\paragraph{Logical dependency boundary.}
The changes in this section preserve the numerical C222 Horn rule predicate
of package 68.  They do not supply a new proof that this predicate is
equivalent to ZFC refutability.  That calibration is inherited from the
repository's pinned C222-to-ZF and constructible-choice arguments, including
the sealed target code $1$ and free truth code $0$.  The equivalence of
complete accumulator certificates with finite Horn proofs is the pinned
single-conclusion flow theorem, composed with the bounded-factor carrier
theorem.  Given those dependencies, all subsequent transformations here are
explicit arithmetic equivalences provable in ZFC; none assumes
$\operatorname{Con}(\mathrm{ZFC})$.

The pertinent sources are pinned to the repository snapshot on the
title page:
\begin{itemize}
\item The DiophantineConsistency Lean modules for bounded-forall
      certificates and their representation, Pell equations,
      exponentiation, and factorials.
\item Package 68's \nolinkurl{COMMUTATIVE_CERTIFICATES.md} and
      \nolinkurl{affine_shared_step.json}.
\item The candidate \nolinkurl{direct-flow-large-carrier/README.md}.
\end{itemize}
The provenance inventory supplies their complete repository paths and
the immediate source copies.

\section{Specializing every power graph to positive arguments}

The completed source can choose $t\ge1$ and $m\ge1$ without losing a
certificate.  A zero-length accumulator path contradicts its distinct
endpoints, and if all labels were zero then every accumulator update would
leave the state unchanged.  Hence the existence predicate is preserved by
the parameterizations $t=t_0+1$, $m=m_0+1$.

Likewise each encoded CRT coordinate may be required positive.  The Chinese
remainder theorem specifies residue classes only; adding a positive multiple
of $B$ preserves all required congruences.  Thus one may parameterize
$A_j=\alpha_j+1$ for natural $\alpha_j$.  This affects encoded coordinates,
not the decoded row witnesses: the latter may still be zero, including the
free truth code.  The range certificate remains valid because its proof
already establishes divisibility of the binomial value for every CRT
representative of the prescribed small residues.

Under these choices every power graph used in the construction has positive
base and positive exponent.  The carrier and modulus powers have exponent
$t>0$; the factorial graphs are called at $t$ and $H$, both positive;
their radius and auxiliary binomial arguments are positive.  In each
coordinate-range binomial graph the upper argument $A_j$ is positive and
the lower argument $u_j+1$ is positive.  Therefore its three power calls
$2^{A_j}$, $(2^{A_j}+2)^{A_j}$, and $(2^{A_j}+1)^{u_j+1}$ are positive
cases.  The zero-base and zero-exponent branches of a general power graph
are unnecessary in this specialized construction.  Their removal is
justified by this argument inventory; it is not a correct optimization of
the unrestricted exponentiation relation.

For the positive-power certificate, its final Pell norm gives
\[
 A^2=1+W^3(W+2)Z^2,
 \qquad W\ge n\ge1,\quad A\ge2.
\]
Consequently $Z\ge1$ and $A>W\ge n$.  Hence the natural subtraction
$A\mathbin{\dot{-}}n$ in the original power theorem equals ordinary
integer subtraction.  The congruence can be written with nonnegative
polynomial sides as
\[
 X+Yn+Tq_+=YA+m+Tq_-.
\]
This is the only subtraction-elimination point that needs a separate
inequality argument in the power specialization.
